\documentclass[10pt,a4paper]{amsart}
\usepackage{amsmath,amssymb,amsthm,graphicx}
\usepackage{verbatim}
\newenvironment{reference}{\comment}{\endcomment}
%\newenvironment{reference}{}{}
\newenvironment{slogan}{\comment}{\endcomment}
\newenvironment{history}{\comment}{\endcomment}

% For commutative diagrams we use Xy-pic
\usepackage[all]{xy}

% We use 2cell for 2-commutative diagrams.
\xyoption{2cell}
% Theorem environments.
%
\theoremstyle{plain}
\newtheorem{theorem}[subsection]{Theorem}
\newtheorem{proposition}[subsection]{Proposition}
\newtheorem{lemma}[subsection]{Lemma}

\theoremstyle{definition}
\newtheorem{definition}[subsection]{Definition}
\newtheorem{example}[subsection]{Example}
\newtheorem{exercise}[subsection]{Exercise}
\newtheorem{situation}[subsection]{Situation}

\theoremstyle{remark}
\newtheorem{remark}[subsection]{Remark}
\newtheorem{remarks}[subsection]{Remarks}

\numberwithin{equation}{subsection}

% Macros
%
\def\lim{\mathop{\mathrm{lim}}\nolimits}
\def\colim{\mathop{\mathrm{colim}}\nolimits}
\def\Spec{\mathop{\mathrm{Spec}}}
\def\Hom{\mathop{\mathrm{Hom}}\nolimits}
\def\Ext{\mathop{\mathrm{Ext}}\nolimits}
\def\SheafHom{\mathop{\mathcal{H}\!\mathit{om}}\nolimits}
\def\SheafExt{\mathop{\mathcal{E}\!\mathit{xt}}\nolimits}
\def\Sch{\mathit{Sch}}
\def\Mor{\mathop{\mathrm{Mor}}\nolimits}
\def\Ob{\mathop{\mathrm{Ob}}\nolimits}
\def\Sh{\mathop{\mathit{Sh}}\nolimits}
\def\NL{\mathop{N\!L}\nolimits}
\def\CH{\mathop{\mathrm{CH}}\nolimits}
\def\proetale{{pro\text{-}\acute{e}tale}}
\def\etale{{\acute{e}tale}}
\def\QCoh{\mathit{QCoh}}
\def\Ker{\mathop{\mathrm{Ker}}}
\def\Im{\mathop{\mathrm{Im}}}
\def\Coker{\mathop{\mathrm{Coker}}}
\def\Coim{\mathop{\mathrm{Coim}}}

% Boxtimes
%
\DeclareMathSymbol{\boxtimes}{\mathbin}{AMSa}{"02}

%
% Macros for moduli stacks/spaces
%
\def\QCohstack{\mathcal{QC}\!\mathit{oh}}
\def\Cohstack{\mathcal{C}\!\mathit{oh}}
\def\Spacesstack{\mathcal{S}\!\mathit{paces}}
\def\Quotfunctor{\mathrm{Quot}}
\def\Hilbfunctor{\mathrm{Hilb}}
\def\Curvesstack{\mathcal{C}\!\mathit{urves}}
\def\Polarizedstack{\mathcal{P}\!\mathit{olarized}}
\def\Complexesstack{\mathcal{C}\!\mathit{omplexes}}
% \Pic is the operator that assigns to X its picard group, usage \Pic(X)
% \Picardstack_{X/B} denotes the Picard stack of X over B
% \Picardfunctor_{X/B} denotes the Picard functor of X over B
\def\Pic{\mathop{\mathrm{Pic}}\nolimits}
\def\Picardstack{\mathcal{P}\!\mathit{ic}}
\def\Picardfunctor{\mathrm{Pic}}
\def\Deformationcategory{\mathcal{D}\!\mathit{ef}}
\usepackage{cleveref}
\crefname{originalreference}{original reference}{original references}
\Crefname{originalreference}{Original reference}{Original references}
\makeatletter
\newlabel{lemma-exact-sequence-NL}{{00S2}{source}{}{sourcepointer-1}{}}
\newlabel{lemma-exact-sequence-NL@cref}{{[originalreference][0][]00S2}{[1][0][]source}{}{}{}}
\makeatother
\crefname{originalreference}{original reference}{original references}
\Crefname{originalreference}{Original reference}{Original references}
\makeatletter
\newlabel{lemma-rig-etale}{{0ALQ}{source}{}{sourcepointer-2}{}}
\newlabel{lemma-rig-etale@cref}{{[originalreference][0][]0ALQ}{[1][0][]source}{}{}{}}
\makeatother
\crefname{originallemma}{original lemma}{original lemmas}
\Crefname{originallemma}{Original lemma}{Original lemmas}
\makeatletter
\newlabel{more-algebra-lemma-henselization}{{0A02}{source}{}{sourcepointer-3}{}}
\newlabel{more-algebra-lemma-henselization@cref}{{[originallemma][0][]0A02}{[1][0][]source}{}{}{}}
\makeatother
\crefname{originallemma}{original lemma}{original lemmas}
\Crefname{originallemma}{Original lemma}{Original lemmas}
\makeatletter
\newlabel{more-algebra-lemma-henselization-Noetherian-pair}{{0AGV}{source}{}{sourcepointer-4}{}}
\newlabel{more-algebra-lemma-henselization-Noetherian-pair@cref}{{[originallemma][0][]0AGV}{[1][0][]source}{}{}{}}
\makeatother
\crefname{originallemma}{original lemma}{original lemmas}
\Crefname{originallemma}{Original lemma}{Original lemmas}
\makeatletter
\newlabel{more-algebra-lemma-quotient-by-idempotent}{{0ALR}{source}{}{sourcepointer-5}{}}
\newlabel{more-algebra-lemma-quotient-by-idempotent@cref}{{[originallemma][0][]0ALR}{[1][0][]source}{}{}{}}
\makeatother
\crefname{originallemma}{original lemma}{original lemmas}
\Crefname{originallemma}{Original lemma}{Original lemmas}
\makeatletter
\newlabel{more-morphisms-lemma-check-smoothness-on-infinitesimal-nbhds}{{0A43}{source}{}{sourcepointer-6}{}}
\newlabel{more-morphisms-lemma-check-smoothness-on-infinitesimal-nbhds@cref}{{[originallemma][0][]0A43}{[1][0][]source}{}{}{}}
\makeatother
\title{Completion maps from etale-over-complement algebras come from the henselization}
\author{The Stacks Project Authors}
\date{stacks-project:restricted.tex:local-snapshot-20261012; GNU FDL 1.2 or later}
\begin{document}
\maketitle
\noindent\textit{Editorial note.} All delivered mathematical text is the author's own, copied from the pinned Stacks Project chapter file restricted.tex; no statement or proof step was rewritten, reordered or removed.\par\medskip
\noindent\textit{Editorial note.} The counted proof is the author's complete argument: the proof is complete: uniqueness, the factorization assertion and existence via the reduction A = C\par\medskip
\noindent\textit{Editorial note.} The label lemma-fully-faithful-etale-over-complement names the counted statement, and the proof environment immediately following it in the source is delivered verbatim as its proof.\par\medskip
\noindent\textit{Editorial note.} The counted statement is carried by the author's own lemma environment in the chapter source.\par\medskip
\noindent\textit{Allowed standard prerequisites (editorial).}\begin{itemize}
\item I-adic completions and the henselization of a Noetherian pair as a filtered colimit of etale algebras preserving the reduction
\item the cotangent complex exact sequence for compositions and the annihilation of its cohomology in the rig-etale case
\item idempotent-quotient lemma identifying a surjective completion with the completion of a finite type algebra
\item etale-ness checked on infinitesimal neighbourhoods
\item the universal property of the henselization
\end{itemize}\medskip
\ifdefined\hypertarget\hypertarget{sourcepointer-1}{}\fi
\noindent\textit{Source reference.} Original reference 00S2: sources/stacks-restricted/pinned-original-source.tex, source lines 725--725. Prerequisite pointer invoked at proof line(s) 2103; it names a standard fact of the same chapter that the counted argument uses, and every step of the counted proof built on it is retained verbatim in the delivered body.\par\medskip
\ifdefined\hypertarget\hypertarget{sourcepointer-2}{}\fi
\noindent\textit{Source reference.} Original reference 0ALQ: sources/stacks-restricted/pinned-original-source.tex, source lines 1981--1981. Prerequisite pointer invoked at proof line(s) 2103; it names a standard fact of the same chapter that the counted argument uses, and every step of the counted proof built on it is retained verbatim in the delivered body.\par\medskip
\ifdefined\hypertarget\hypertarget{sourcepointer-3}{}\fi
\noindent\textit{Source reference.} Original lemma 0A02: sources/stacks-chow/tags.txt, source lines 12620--12620. Prerequisite pointer invoked at proof line(s) 2076; it names a standard fact of a sibling chapter that the counted argument uses, and every step of the counted proof built on it is retained verbatim in the delivered body.\par\medskip
\ifdefined\hypertarget\hypertarget{sourcepointer-4}{}\fi
\noindent\textit{Source reference.} Original lemma 0AGV: sources/stacks-chow/tags.txt, source lines 13208--13208. Prerequisite pointer invoked at proof line(s) 2084; it names a standard fact of a sibling chapter that the counted argument uses, and every step of the counted proof built on it is retained verbatim in the delivered body.\par\medskip
\ifdefined\hypertarget\hypertarget{sourcepointer-5}{}\fi
\noindent\textit{Source reference.} Original lemma 0ALR: sources/stacks-chow/tags.txt, source lines 13379--13379. Prerequisite pointer invoked at proof line(s) 2110; it names a standard fact of a sibling chapter that the counted argument uses, and every step of the counted proof built on it is retained verbatim in the delivered body.\par\medskip
\ifdefined\hypertarget\hypertarget{sourcepointer-6}{}\fi
\noindent\textit{Source reference.} Original lemma 0A43: sources/stacks-chow/tags.txt, source lines 12761--12761. Prerequisite pointer invoked at proof line(s) 2119; it names a standard fact of a sibling chapter that the counted argument uses, and every step of the counted proof built on it is retained verbatim in the delivered body.\par\medskip
\begin{lemma}
\label{lemma-fully-faithful-etale-over-complement}
Let $A$ be a Noetherian ring. Let $I \subset A$ be an ideal.
Let $B$ be a finite type $A$-algebra such that
$\Spec(B) \to \Spec(A)$ is \'etale over $\Spec(A) \setminus V(I)$.
Let $C$ be a Noetherian $A$-algebra. Then any $A$-algebra
map $B^\wedge \to C^\wedge$ of $I$-adic completions
comes from a unique $A$-algebra map
$$
B \longrightarrow C^h
$$
where $C^h$ is the henselization of the pair $(C, IC)$ as
in More on Algebra, Lemma \ref{more-algebra-lemma-henselization}.
Moreover, any $A$-algebra homomorphism $B \to C^h$ factors through
some \'etale $C$-algebra $C'$ such that $C/IC \to C'/IC'$ is an isomorphism.
\end{lemma}

\begin{proof}
Uniqueness follows from the fact that $C^h$ is a subring of
$C^\wedge$, see for example
More on Algebra, Lemma \ref{more-algebra-lemma-henselization-Noetherian-pair}.
The final assertion follows from the fact that $C^h$ is the filtered colimit
of these $C$-algebras $C'$, see proof of
More on Algebra, Lemma \ref{more-algebra-lemma-henselization}.
Having said this we now turn to the proof of existence.

\medskip\noindent
Let $\varphi : B^\wedge \to C^\wedge$ be the given map.
This defines a section
$$
\sigma : (B \otimes_A C)^\wedge \longrightarrow C^\wedge
$$
of the completion of the map $C \to B \otimes_A C$. We may
replace $(A, I, B, C, \varphi)$ by $(C, IC, B \otimes_A C, C, \sigma)$.
In this way we see that we may assume that $A = C$.

\medskip\noindent
Proof of existence in the case $A = C$. In this case the map
$\varphi : B^\wedge \to A^\wedge$ is necessarily surjective.
By Lemmas \ref{lemma-rig-etale} and \ref{lemma-exact-sequence-NL}
we see that the cohomology groups of
$\NL_{A^\wedge/\!_\varphi B^\wedge}^\wedge$
are annihilated by a power of $I$. Since $\varphi$ is surjective,
this implies that $\Ker(\varphi)/\Ker(\varphi)^2$ is annihilated
by a power of $I$. Hence $\varphi : B^\wedge \to A^\wedge$
is the completion of a finite type $B$-algebra $B \to D$, see
More on Algebra, Lemma \ref{more-algebra-lemma-quotient-by-idempotent}.
Hence $A \to D$ is a finite type algebra map which induces an isomorphism
$A^\wedge \to D^\wedge$. By
Lemma \ref{lemma-rig-etale} we may replace $D$ by a localization
and assume that $A \to D$ is \'etale away from $V(I)$.
Since $A^\wedge \to D^\wedge$ is an isomorphism, we see that
$\Spec(D) \to \Spec(A)$ is also \'etale in a neighbourhood of $V(ID)$
(for example by
More on Morphisms, Lemma
\ref{more-morphisms-lemma-check-smoothness-on-infinitesimal-nbhds}).
Thus $\Spec(D) \to \Spec(A)$ is \'etale. Therefore $D$ maps to
$A^h$ and the lemma is proved.
\end{proof}

\end{document}
